\documentclass[11pt]{article}
\usepackage{paperreview}
\title{When Memory Is Observable:\\Control Bounds and Temporal-Coarse-Graining Obstructions}
\author{Luqman Mahmood}
\date{1 October 2026}
\hypersetup{pdftitle={When Memory Is Observable: Control Bounds and Temporal-Coarse-Graining Obstructions},pdfauthor={Luqman Mahmood}}
\begin{document}
\maketitle
\draftstatus{SM2, v1.0.0}{Conditional bounds and interval-certified scales; no performed physical experiment.}
\begin{abstract}
An increase of distinguishability in an averaged qubit evolution can be
dynamically real yet difficult to identify from measured records. For a qubit
coupled along a fixed axis to a symmetric classical telegraph process, we
derive, in the fast-switching regime $\gamma\geq4v$ ($0<v/\gamma\leq1/4$),
an all-control ceiling of order $v^2D^2/\gamma^4$ on each endpoint
increase of squared trace distance when transverse drive strength is bounded.
An explicit constant-drive construction attains that order. We then consider
unresolved timing: without a common bound on control speed, one continuous
unitary trajectory can approximate the entire time-averaged path of any
continuous unital qubit channel. This gives a uniform finite-sample
indistinguishability result for independently reset, single-time experiments.
A separate, externally calibrated set of speed-and-clock assumptions restores a
nonempty distinguishability region. For a specified eight-setting,
point-time experiment, information bounds make a fourth-power trial cost
necessary. Independent readout flips at a fixed nonzero rate raise this
necessary order to the eighth power, even with exactly known mean bias.
Standard quantum-channel theory and
classical noise are assumed. No physical bath, instrument precision,
experimental confidence or practicable shot count is demonstrated.
\end{abstract}

\section{A dynamics claim is not yet a measurement claim}
The relevant comparison is between two explanations for an observed
change in distinguishability: an averaged evolution with memory, and an
ordinary fine-time evolution viewed through an unresolved observation clock.
The comparison class, reset rule and calibration are parts of the theorem.
Trace-distance backflow is the established Breuer--Laine--Piilo witness
\cite{BLP}; completely positive divisibility and correlation-based criteria
are distinct notions \cite{RHP,BLPV}. This paper bounds individual returns
of squared trace distance, not the complete positive-variation measure,
and does not identify that witness with full multitime process memory.
Classical telegraph and colored noise already produce memory effects
\cite{Daffer,Benedetti2013,Benedetti2014}. The contribution sought here is
the bounded-drive fast-switching estimate and its acquisition constraints,
not memory from classical noise itself.

Unital-qubit canonical forms and decompositions into unitary channels are
established \cite{RSW,LiChoi}. We use them rather than claim a new channel
classification. The candidate contributions are an explicit bounded-drive
visibility law, a time-path construction with a uniform error estimate, and
their operational separation. Comparison with the broader control,
relaxation and statistical-testing literature has not been attempted here;
no priority is claimed for these bounds.
Exact telegraph-noise dynamics and pulse control are established
\cite{BergliFaoro}; the present drive contract instead excludes
instantaneous pulses of unbounded transverse strength. The time-path
construction is a quantitative application of relaxed-control/chattering
ideas \cite{Artstein}, not a claim to have discovered rapid-control
approximation. Its observation model is independently reset single-time
experiments, not a general multitime quantum-process tomography theorem.

\section{Classical telegraph model and conventions}
Let a symmetric two-state process switch at rate $\gamma>0$, and couple to
a qubit with Hamiltonian contribution $\xi(t)v\sigma_z$, $\xi=\pm1$.
A prescribed system Hamiltonian is $b(t)\cdot\sigma$.
For half the Bloch-vector difference of any two preparations write
$m$ for the mean and $h$ for the signed conditional difference.
In time $s=\gamma t$,
\begin{equation}\label{eq:tele}
 m'=\Omega m+2\epsilon Jh,\qquad
 h'=\Omega h+2\epsilon Jm-2h,
 \quad Jx=e_z\times x,\quad\epsilon=v/\gamma .
\end{equation}
The independent balanced initial noise state gives $h(0)=0$ and
$\norm{m(0)}\leq1$. Require
\[
 \norm{[\Omega(s),J]}\leq2\delta,\qquad
 \delta=D/\gamma ,
\]
equivalently $\sqrt{b_x^2+b_y^2}\leq D$ in the physical Hamiltonian.
The longitudinal field is locally bounded but need not obey the same cap.
Controls are piecewise continuous and prescribed, not feedback using the
hidden fluctuator. Neither a bandwidth cap nor an instantaneous-pulse model
is implied.

The trace distance of the two mean states is $\norm m$.
Skew symmetry gives the exact balance
\begin{equation}\label{eq:balance}
 (\norm m^2+\norm h^2)'=-4\norm h^2.
\end{equation}
This is a distinguishability identity, not a thermodynamic-energy law.

\section{A uniform endpoint-return ceiling}
\begin{theorem}\label{thm:ceiling}
For $0<\epsilon\leq1/4$ and every admitted control and preparation,
\[
 \norm{m(s_2)}^2-\norm{m(s_1)}^2
 \leq\frac{16\epsilon^2\delta^2}
 {(1+\sqrt{1-4\epsilon^2})^4}
 <2\epsilon^2\delta^2
 \quad(s_2\geq s_1).
\]
The final strict inequality is non-strict in a degenerate zero-coupling
or zero-transverse-drive case.
\end{theorem}
\begin{proof}
Put $\beta=\sqrt{1-4\epsilon^2}$, $\lambda=1+\beta$ and
\[
 k(s)=
 \frac{2\epsilon(1-\e^{-2\beta s})}
 {1+\beta-(1-\beta)\e^{-2\beta s}} .
\]
Then $k'=2\epsilon-2k+2\epsilon k^2$, $k(0)=0$,
and $0\leq k\leq2\epsilon/\lambda$.
The residual $r=h-kJm$ obeys
\[
 r'=(\Omega-2I-2\epsilon kJ^2)r+k[\Omega,J]m .
\]
Its homogeneous symmetric part is at most $-\lambda I$.
The balance \eqref{eq:balance} gives $\norm m\leq1$, hence
\[
 \norm r\leq2\delta\int_0^s\e^{-\lambda(s-u)}k(u)\dd u
 \leq4\epsilon\delta/\lambda^2=:R.
\]
Let $f=\norm m^2$, $g=\norm h^2$ and $x=Jm$. Direct calculation gives
$f'=-4\epsilon(k\norm x^2+x\cdot r)$.
Whenever $f'>0$,
\[
 g=\norm r^2+2k(k\norm x^2+x\cdot r)-k^2\norm x^2\leq R^2.
\]
If a positive derivative occurs between two endpoints, let $\tau$ be the
infimum of such times. Before $\tau$, $f$ is nonincreasing; continuity
gives $g(\tau)\leq R^2$. Nonincrease of $f+g$ then gives
$f(s_2)-f(s_1)\leq g(\tau)-g(s_2)\leq R^2$.
If there is no positive derivative, the conclusion is immediate.
Finally $\lambda>7/4$ and $16/(7/4)^4<2$.
\end{proof}
The theorem bounds each endpoint increase. It does not bound the sum of
all positive variations over an arbitrarily long record.
The fast-switching assumption is essential to the stated ceiling. With no
drive and an initially transverse vector, eliminating $h$ gives
$m''+2m'+4\epsilon^2m=0$. For $\epsilon>1/2$, its scalar factor is
$\e^{-s}(\cos\omega s+\sin\omega s/\omega)$,
$\omega=\sqrt{4\epsilon^2-1}$, and it has distinguishability revivals even
though $D=0$. The present ceiling makes no claim in that regime.

\section{A matching-order construction}
\begin{proposition}\label{prop:lower}
For $0<\epsilon,\delta\leq1/32$, use a constant drive $D\sigma_x$ and
\[
 m(0)=\frac{(0,2\delta,1-\delta^2)}{1+\delta^2}.
\]
Between $s=1$ and $s=5/4$ the squared-distance return $\Delta$ satisfies
\[
 \frac{\epsilon^2\delta^2}{32}<\Delta<2\epsilon^2\delta^2.
\]
\end{proposition}
\begin{proof}
The invariant subsystem for
$H=h_x/(\epsilon\delta)$, $Y=m_y/\delta$, $Z=m_z$ is
\[
 H'=-2H-2Y,\quad Y'=2\epsilon^2H-2Z,\quad Z'=2\delta^2Y.
\]
The rescaled zero-parameter reference is
$Z_0=1$, $Y_0=2-2s$, $H_0=2s-3+3\e^{-2s}$.
The exact return is $4\epsilon^2\delta^2\int_1^{5/4}YH\dd s$,
whose reference coefficient is
\[
 c_0=\frac16-6\e^{-2}+9\e^{-5/2}>\frac9{100}.
\]
The integral equations on $[0,5/4]$ give
$\sup|Y|<5$, $|Z-1|\leq(29/2)\delta^2$, and
\[
 |Y-Y_0|,\ |H-H_0|\leq
 E:=\frac{25}{2}\epsilon^2+\frac{153}{4}\delta^2\leq\frac{203}{4096}.
\]
On the integration interval, $|Y_0|\leq1/2$ and $|H_0|<3/5$.
Consequently the coefficient error is at most $11E/10+E^2$.
The rational inequality
$9/100-11(203/4096)/10-(203/4096)^2>1/32$ proves the lower bound.
The upper bound is Theorem \ref{thm:ceiling}.
\end{proof}
At fixed $v,D$, the order is $\gamma^{-4}$. The preparation depends on
$D/\gamma$ and the observation window scales as $1/\gamma$.
This is not a fixed-preparation, fixed-clock prediction.

\section{A whole-path timing obstruction}
Write
$d(\Phi,\Psi)=\frac12\norm{\Phi-\Psi}_\diamond$.
It is at most one for channels and bounds all single-outcome probability
differences, including ancilla-assisted ones.

Every unital qubit channel is an equal average of four unitary channels
\cite{LiChoi}. For a diagonal canonical channel with Bell--Choi
probabilities $p_0,\ldots,p_3$, one construction is
\[
 U_k=\sqrt{p_0}I-\ii\sum_{a=1}^3s_{ka}\sqrt{p_a}\sigma_a ,
\]
where the sign rows are $(+++),(+--),(-+-),(--+)$. Averaging kills
the cross terms. Outer unitary conjugations restore the general channel.

\begin{theorem}[Time-averaged closure]\label{thm:clock}
Let $\mathcal C(s)$ be a diamond-continuous unital-qubit channel path on
$[0,T]$, with $0<w<T/2$, and define
\[
 \overline{\mathcal C}(t)=\frac1{2w}
 \int_{t-w}^{t+w}\mathcal C(s)\dd s,\quad w\leq t\leq T-w.
\]
For every $\varepsilon>0$ there is one continuous, piecewise-smooth
$U(s)\in SU(2)$ with $U(0)=I$ whose averaged conjugation channel is
uniformly within $\varepsilon$ of $\overline{\mathcal C}$.
Each chosen path has finite-speed pieces, but no common speed cap as
$\varepsilon\downarrow0$ is asserted.
\end{theorem}
\begin{proof}
Divide time into equal cells of length $\ell$. Freeze the target at each
left endpoint and dwell equally at its four representing unitaries.
If $\omega(\ell)$ is the target modulus of continuity, freezing costs
at most $\omega(\ell)$. Full cells have the exact same average; the
two possible partial cells cost at most $\ell/w$.

Replace each jump by an $SU(2)$ arc in the first $\tau\leq\ell/4$ of the
next subcell, including an initial arc from $I$.
At most $8w/\ell+2$ arc intervals intersect an averaging window.
Their changed fraction is at most $4\tau/\ell+\tau/w$. Therefore
\[
 \sup_t d(\overline{\mathcal U}(t),\overline{\mathcal C}(t))
 \leq\omega(\ell)+\ell/w+4\tau/\ell+\tau/w .
\]
Choose $\omega(\ell)\leq\varepsilon/4$, $\ell\leq\varepsilon w/4$ and
$\tau\leq\varepsilon\ell/16$, reducing $\varepsilon$ to at most one if
needed. The bound is at most $49\varepsilon/64$.
The connecting arcs have finite generators for each positive $\tau$.
\end{proof}
This approximates an \emph{already averaged} target path. It does not
realize arbitrary incompatible maps at overlapping windows, nor does it
cover nonunital channels or higher-dimensional quantum Birkhoff problems.

\begin{corollary}[Uniform finite-sample obstruction]
Assume independent reset trials and fresh independent uniform clock draws;
only nominal times and outcomes are recorded. For any finite number $N$ of
single-time channel uses, a test with size at most $\alpha$ for every
admitted fine-time unitary path without a common speed cap has rejection
probability at most $\alpha$ on the averaged target.
\end{corollary}
\begin{proof}
The uniform channel approximation gives a record-law total-variation
distance at most $N\varepsilon$ by the hybrid argument, including adaptive
input/effect choices and single-time ancillas. The rejection probability is
at most $\alpha+N\varepsilon$ for every $\varepsilon>0$.
\end{proof}
The approximating path may change with $\varepsilon$. This is not one fixed
finite-speed null with identical laws at arbitrarily large sample sizes.
Within-trial multitime interventions and time-resolved records are outside
the experiment.

\section{What external calibration restores}
Under an independently justified pair-state speed cap of $4$, the special
target $v=D=1$, $\gamma\geq32$ admits a sufficient clock-and-error region
\[
 w\leq\frac1{256\gamma^2},\qquad
 \eta_{\rm other}\leq\frac{\gamma^{-4}}{2400}.
\]
The squared distinguishability witness from three first-time orthogonal
axes and one free-rotated final axis has an all-centres ideal margin
greater than $\gamma^{-4}/500$. The null clock allowance is
$16w^2/3\leq\gamma^{-4}/12288$, giving a clock-aware margin greater than
$\gamma^{-4}/600$. The target's probability curvature bounds its centered
jitter error by $7w^2/6$; that error and $\eta_{\rm other}$ stay below
the required total $\gamma^{-4}/1200$ budget.
These inequalities are a conditional acquisition region, not achieved
instrument performance. Their detailed calibration proof is included in
the supplement as the bounded-drive and tomography-clock proofs.
They use a different set of assumptions from
the unrestricted-speed obstruction.

The operational prediction, within $\gamma\geq4v$, is that increasing the
switching rate at fixed coupling and drive makes the best endpoint visibility scale no better than
the fourth inverse power, while a tailored admitted preparation attains
that order. More samples cannot compensate for an unrestricted clock/control
null class; an external bound must first narrow that class.

\section{Statistical cost and the readout trade-off}
\label{sec:measurement-cost}
We now distinguish a probability-error allowance from the number of trials
required to attain it. Throughout this section $v=D=1$, $\gamma\geq32$,
$\delta=\gamma^{-1}$ and $\sigma=\gamma^{-4}$. At $s=1$ measure both
opposed preparations in the orthonormal triad
\[
 f_1,\quad e_x,\quad (0,f_{1z},-f_{1y}),
\]
and at $s=5/4$ along $f_2$, where $f_i$ is the freely rotated initial
vector at the corresponding time. All axes and counts are fixed before
acquisition. This gives eight binary settings, with pair gains
$g_1,g_2,g_3,g_4$ and witness
$W=g_4^2-g_1^2-g_2^2-g_3^2$.

The following control and axis-error estimates for the target are proved
in the supplement \cite{Internal}:
\begin{gather}
 g_2=0,\qquad \sigma/48<W<2\sigma,\qquad
 q_1,q_4\leq C_0\sigma,\quad C_0=7109689/4194304<7/4,\label{eq:target-bounds}\\
 |g_3|\leq(85/1024)\sqrt\sigma,\qquad
 \norm{m_1-f_1}\leq A\delta^3,\quad A=85/32.
\end{gather}
Here $q_i=(1-g_i)/2$ are the aligned failures for either preparation of
this unital target. The detector and the null do not assume unitality.
These estimates are not consequences of a finite numerical scan.
Their proofs are in Section~6 of \emph{A bounded-strength endpoint-return
law} and Section~7 of \emph{Time-averaged tomography, closure, and an
admissible region}. The separate \emph{Finite-sample rare-event return
witness} takes these as inputs for the count argument. The supplement
README maps these titles to the supplied files.
Finite interval controls do not replace those proofs.

\subsection{A sound finite-count rule}
Let $a=k/n$ be the observed failure fraction in an aligned setting, and
let $\vartheta>0$ be the confidence parameter. The data-derived radius
\[
 R(a,n,\vartheta)=\sqrt{2a\vartheta/n}+4\vartheta/n
\]
covers the true Bernoulli mean with error probability at most
$2e^{-\vartheta}$, for every mean, including nonrare nulls. Indeed Bernstein's
bound gives $|a-q|\leq\sqrt{2q\vartheta/n}+2\vartheta/(3n)$.
For $q>a$, solving the quadratic gives
\[
 q-a\leq\frac{5\vartheta}{3n}
       +\sqrt{\frac{2a\vartheta}{n}+\frac{7\vartheta^2}{3n^2}}
       <R(a,n,\vartheta);
\]
for $q\leq a$ the result is immediate. In particular, zero observed
failures still leave radius $4\vartheta/n$.

At an aligned axis use $\widehat g=1-a_+-a_-$ and
$e=R_++R_-+2\beta$, where $\beta$ bounds each probability's nonclock
systematic error. At a transverse axis use the difference of positive-outcome
frequencies and $e=\sqrt{\vartheta/(2n_+)}+
\sqrt{\vartheta/(2n_-)}+2\beta$. Reject only if
\begin{equation}\label{eq:count-rule}
 \max(0,|\widehat g_4|-e_4)^2
 -\sum_{j=1}^3\min(1,|\widehat g_j|+e_j)^2>16w^2/3.
\end{equation}
IID reset trials within each setting and the prescribed clock law imply
simultaneous coverage except with probability $16e^{-\vartheta}$.
On coverage, the expression is a lower bound on the clock-averaged witness;
the qualified contractive null cannot exceed $16w^2/3$.
Soundness uses no rare-event property of the null.

With $w\leq1/(256\gamma^2)$ and $\beta\leq\sigma/100000$, the detailed
target power estimate in the supplied confidence/count proof gives the sufficient count
\begin{equation}\label{eq:sufficient-count}
 N_{\rm suff}=4\lceil2^{22}\vartheta\gamma^4\rceil
             +4\lceil2^{16}\vartheta\gamma^4\rceil
\end{equation}
for both errors at most $16e^{-\vartheta}$, $\vartheta\geq3$.
The first four counts use aligned failures; the other four use transverse
positive outcomes. This is a conservative sufficient guarantee, not an
optimal test. It includes both displacement of the empirical centre and
the confidence radius, and does not divide systematic error by sample size.

\subsection{Why fourth-power cost cannot be removed at point times}
The following lower bound applies to independent single-time trials chosen
from these eight settings at $w=0$. It is not a bound on arbitrary multitime
quantum experiments. The null includes CP-divisible rotations and isotropic
depolarization, with pair-state speed at most four in physical time.

\begin{theorem}[Point-time information obstruction]\label{thm:information}
If both errors of such a test are at most $16e^{-\vartheta}$, where
$\vartheta>\log64$, then its total trial count satisfies
\[
 N>2(\vartheta-\log64)(1-6\sigma)\gamma^4.
\]
\end{theorem}
\begin{proof}
Put $C=7/4$, $b_*=1-2C\sigma>1/2$ and
$r_1=\sqrt{g_1^2+g_3^2}\geq g_1\geq b_*$.
A null can match the first-time vector, then rotate at fixed norm $r_1$
to the final measurement axis. Its only differing setting probabilities
have shift
\[
 d=(g_4-r_1)/2=\frac{W}{2(g_4+r_1)}<\frac{\sigma}{2b_*}.
\]
To check admissibility, interpolate the norm from one to $r_1$ by isotropic
depolarization and rotate along the shorter arc from the initial vector to
$m_1/r_1$; then rotate at fixed norm to $f_2$.
Here $b_*^2\leq1-A^2\delta^6$, so the vector-error bound and
$\arcsin u\leq u/\sqrt{1-u^2}$ give angle bounds
$2\delta+A\delta^3/b_*$ and $\delta/2+A\delta^3/b_*$.
Linear schedules have pair speeds at most
$2+A\delta^2/b_*+2C\delta^3$ and $2+4A\delta^2/b_*$, both below three.
The path is therefore an admitted CP-divisible null, not merely a list of
unrelated endpoints.

For a changed failure probability $q$ and $q'=q+d$, the elementary
Bernoulli relative-entropy bound gives
\begin{align*}
 D_{\rm KL}(q\Vert q')&\leq\frac{d^2}{q'(1-q')}
 \leq\frac d{1-q'}\\
 &<\frac{\sigma}{2\{b_*(1-C\sigma)-\sigma/2\}}
 \leq\frac{\sigma}{2(1-6\sigma)}.
\end{align*}
The final inequality uses $3C+1/2=23/4<6$.
The Bretagnolle--Huber testing inequality gives
$\alpha+\beta_{\rm II}\geq\tfrac12\exp(-D_{\rm KL}(P\Vert Q))$
\cite[Cor.~4]{Canonne}. Since the total information is at most $N$ times
the largest per-trial divergence and $\alpha+\beta_{\rm II}\leq
32e^{-\vartheta}$, the stated result follows.
\end{proof}
The point-time case is a necessary obstruction for any rule uniform over
an allowed clock class containing $w=0$. Constants for a fixed nonzero
jitter require the actual averaged target and null probabilities.

\subsection{Systematic masking and random detector errors are different}
Within a split-null family that increases $g_1$ by $2d$ and decreases
$g_4$ by $2d$, equal endpoint probability shifts satisfy
\[
 (g_4-2d)^2=(g_1+2d)^2+g_3^2,\qquad
 \beta_{\rm split}=\frac d2=\frac{W}{8(g_4+g_1)}.
\]
An error allowance $\beta_{\rm split}$ on each of the two hypotheses
admits their common midpoint law. This is a masking example for
settingwise adversarial bias, not the global minimax threshold.

If all eight settings instead share one affine readout
$p_{\rm obs}=a_{\rm det}+b_{\rm det}p$ with $b_{\rm det}>0$, differences of
opposed preparations cancel the offset exactly:
\[
 g_{j,\rm obs}=b_{\rm det}g_j,\qquad
 W_{\rm obs}=b_{\rm det}^2W.
\]
This preserves the point-time sign. It neither proves shared gain in an
apparatus nor cancels sampling variance.
In particular, symmetric independent outcome flips are this same affine
readout with $a_{\rm det}=r$ and $b_{\rm det}=1-2r$: cancellation of the
offset in the mean does not cancel the Bernoulli variance floor below.

\begin{proposition}[Fixed random readout flips change the necessary order]
Under the hypotheses of Theorem~\ref{thm:information}, suppose each outcome
is independently flipped with a known fixed probability $r\in(0,1/2)$.
Writing $\mu=1-2r$, every such test requires
\[
 N>\frac{4b_*^2r(1-r)(\vartheta-\log64)}{\mu^2}\,\gamma^8.
\]
\end{proposition}
\begin{proof}
An ideal failure probability becomes $r+\mu q$, bounded between $r$ and
$1-r$. The same null changes it by $\mu d$, so
\[
 D_{\rm KL,obs}\leq\frac{\mu^2d^2}{r(1-r)}
       <\frac{\mu^2\sigma^2}{4b_*^2r(1-r)}.
\]
Apply the same testing inequality. Knowing $r$ exactly does not remove
the random flips or their variance floor.
\end{proof}
This proposition is detector-specific. Heralded errors, correlated reference
readings, or a different acquisition are not excluded by it. With ideal
readout the variance per aligned trial is $O(\gamma^{-4})$, whereas the
standard deviation is $O(\gamma^{-2})$; those quantities must not be confused.

\begin{corollary}[Necessary readout-noise crossover]\label{cor:readout-crossover}
Keep the point-time experiment of Theorem~\ref{thm:information} and fix
$\vartheta>\log64$. At each $\gamma\geq32$, suppose all outcomes are
independently flipped with the same known probability
$0<r_\gamma<1/2$. Set $H=\vartheta-\log64$ and
$b_*=1-7\gamma^{-4}/2$. Then every such test whose two error
probabilities are at most $16e^{-\vartheta}$ has total trial count
\[
 N>H\max\left\{2(1-6\gamma^{-4})\gamma^4,\quad
 4b_*^2\frac{r_\gamma(1-r_\gamma)}{(1-2r_\gamma)^2}\gamma^8\right\}.
\]
If $r_\gamma\asymp\gamma^{-a}$ with $a\geq0$ and
$r_\gamma\leq1/2-\eta$ for some fixed $\eta>0$, this gives
$N=\Omega(\gamma^{\max\{4,8-a\}})$. A trial budget
$N=O(\gamma^4)$ therefore requires $r_\gamma=O(\gamma^{-4})$
in this model; this condition alone is not a sufficient protocol.
\end{corollary}
\begin{proof}
The preceding KL inequality is finite-parameter and holds for every
$0<r<1/2$, so it applies directly at $r=r_\gamma$. Data processing under
the independent flip channel also bounds the observed per-trial KL by
the clean per-trial KL in Theorem~\ref{thm:information}. Combining the
two testing bounds gives the maximum. For $\gamma\geq32$, the factors
$2(1-6\gamma^{-4})$ and $4b_*^2$ lie between one and four, independently
of $r_\gamma$. Away from $1/2$, the ratio
$r(1-r)/(1-2r)^2$ is comparable to $r$, proving the stated powers.
Moreover that ratio is at least $r$ on $[0,1/2)$, so the second bound
forces $r_\gamma=O(\gamma^{-4})$ whenever $N=O(\gamma^4)$.
\end{proof}
At $r=0$ the clean bound applies. At $r=1/2$ every observed outcome is a
fair coin under both hypotheses, so no finite number of these observations
meets the stated error requirement. Rates approaching $1/2$ can impose
a stronger obstruction than the power comparison above. The corollary
does not establish a matching upper bound or extend to correlated or
heralded errors, multitime instruments, or a fixed nonzero exposure.

\paragraph{Numerical scale, not an experimental forecast.}
For $\gamma=32$, $\vartheta=10$, a separately coded six-state generator
and split-null calculation at 80 digits give a clean-readout lower bound
of approximately $5.798\times10^{10}$ trials and
$\beta_{\rm split}\simeq5.446\times10^{-9}$. Equation~\eqref{eq:sufficient-count}
gives $178{,}670{,}639{,}513{,}600$ trials. With independent flips at
$r=0.001$, the same numerical null family gives a lower bound near
$9.883\times10^{13}$. These lower bounds are neither exact optima nor
achieved tests. The portable Arb certificate independently encloses these
rounded values using both the six-dimensional $(m,h)$ and the reduced
three-dimensional systems at 256 and 384 bits. It certifies finite controls
at $\gamma=32,64,128,256$, and brackets the best lower bound within the
specified split-null family at the quoted readout settings. It is not a
global minimax optimization or a proof for all $\gamma$ by sampling.
The separate 80-digit calculation remains a floating-point cross-check.
No apparatus, shot rate or physical unit
conversion is supplied. A viable instrument must address stable common
gain \emph{and} sufficiently rare or independently flagged readout errors.


\section{Limitations and next physical test}
A physical test must separately qualify state preparation, independent reset,
the coupling axis, drive strength, clock distribution, gain/offset and
systematic error. The derived trial bounds concern the stated model and
detector contracts, not an achieved shot budget or hardware rate.
The present construction does not identify a quantum bath. The mathematical
obstruction and positive region have different assumptions; neither should
be substituted for the other after observing an outcome.

\aistatement
\begin{thebibliography}{99}
\bibitem{BLP} H.-P. Breuer, E.-M. Laine and J. Piilo,
\emph{Measure for the Degree of Non-Markovian Behavior of Quantum Processes
in Open Systems}, Phys. Rev. Lett. \textbf{103}, 210401 (2009).
\url{https://doi.org/10.1103/PhysRevLett.103.210401}.
\bibitem{RHP} A. Rivas, S. F. Huelga and M. B. Plenio,
\emph{Entanglement and Non-Markovianity of Quantum Evolutions},
Phys. Rev. Lett. \textbf{105}, 050403 (2010).
\url{https://doi.org/10.1103/PhysRevLett.105.050403}.
\bibitem{BLPV} H.-P. Breuer, E.-M. Laine, J. Piilo and B. Vacchini,
\emph{Colloquium: Non-Markovian dynamics in open quantum systems},
Rev. Mod. Phys. \textbf{88}, 021002 (2016).
\url{https://doi.org/10.1103/RevModPhys.88.021002}.
\bibitem{Daffer} S. Daffer, K. W\'odkiewicz, J. D. Cresser and J. K. McIver,
\emph{Depolarizing channel as a completely positive map with memory},
Phys. Rev. A \textbf{70}, 010304(R) (2004).
\url{https://doi.org/10.1103/PhysRevA.70.010304}.
\bibitem{Benedetti2013} C. Benedetti, F. Buscemi, P. Bordone and M. G. A. Paris,
\emph{Dynamics of quantum correlations in colored-noise environments},
Phys. Rev. A \textbf{87}, 052328 (2013).
\url{https://doi.org/10.1103/PhysRevA.87.052328}.
\bibitem{Benedetti2014} C. Benedetti, M. G. A. Paris and S. Maniscalco,
\emph{Non-Markovianity of colored noisy channels},
Phys. Rev. A \textbf{89}, 012114 (2014).
\url{https://doi.org/10.1103/PhysRevA.89.012114}.
\bibitem{RSW} M. B. Ruskai, S. Szarek and E. Werner,
\emph{An Analysis of Completely-Positive Trace-Preserving Maps on $2\times2$
Matrices}, Linear Algebra Appl. \textbf{347}, 159--187 (2002).
\url{https://arxiv.org/abs/quant-ph/0101003}.
\bibitem{LiChoi} C.-K. Li and M.-D. Choi,
\emph{On unital qubit channels} (2023).
\url{https://arxiv.org/abs/2301.01358}.
\bibitem{Canonne} C. L. Canonne,
\emph{A short note on an inequality between KL and TV},
arXiv:2202.07198v2 (2023), Corollary 4.
\url{https://arxiv.org/abs/2202.07198}.
\bibitem{BergliFaoro} J. Bergli and L. Faoro,
\emph{Exact solution for the dynamical decoupling of a qubit with telegraph noise},
Phys. Rev. B \textbf{75}, 054515 (2007).
\url{https://doi.org/10.1103/PhysRevB.75.054515}.
\bibitem{Artstein} Z. Artstein,
\emph{Rapid Oscillations, Chattering Systems, and Relaxed Controls},
SIAM J. Control Optim. \textbf{27}, 940--948 (1989).
\url{https://doi.org/10.1137/0327050}.
\bibitem{Internal} L. Mahmood, Reproducibility supplement to this paper (2026):
calibration and count proofs, independently coded floating-point check,
and Arb interval certificate. See \texttt{supplement/README.md}.
\end{thebibliography}
\end{document}
